\section{Limitations and conclusion}
\label{sec:conclusion}

\paragraph{Limitations.} (i) \emph{Completeness.} Gaia misses most brown dwarfs
and some bright stars. We add the most important ones by hand and bound the
remaining effect by injection, which lowers the threshold by about
\PctBDTwenty\%. The true neighbourhood is therefore somewhat \emph{easier} to
cross than our graph. (ii) \emph{Kinematics.} Stars move on straight lines.
Galactic curvature, binary orbits and close-encounter scattering are ignored.
This is adequate for $\lesssim1$~Myr inside 50~pc but not for Galaxy-scale
extrapolation. (iii) \emph{Missing RVs} are drawn from a conditional velocity
ellipsoid. Their individual values are uncertain, though the statistics are
unbiased. (iv) \emph{Boundaries.} The moving-star results lose inflowing stars
after about 500~kyr and are conservative there. (v) \emph{Engineering} models
are deliberately simple: one range per ship, refuelling as a binary property of
a star, and a sail-braking law scaled from one calibrated case. (vi) We treat
a colony's ability to relaunch as given, with an optional settling delay.
Sociology is out of scope.

\paragraph{Conclusion.} Routed exactly on the measured solar neighbourhood, the
interstellar graph has a sharp percolation threshold near \RFifty~pc per hop,
about twice the typical spacing. That threshold falls by half once ships may
wait $\sim$$10^5$ years for the stars themselves to close the gaps. In this
regime the expansion front moves at stellar speeds. A fusion ship that survives
a couple of centuries per hop crosses the neighbourhood. Sails need luminous
stopping points, and refuelling only at known planet hosts multiplies the
required tank by \TankRatio. The time-dependent formulation is exact, cheap and
scales directly to the 100-pc GCNS and to deeper Gaia releases. Natural next
steps are Galactic-potential orbits, gravitational-assist edges between stars,
and stochastic colony survival.

\paragraph{Reproducibility.} \texttt{make test \&\& make all} in
\texttt{projects/18\_interstellar\_routes} downloads and caches the catalogue,
runs every experiment, and regenerates the figures and every number in this
paper (about 3~minutes on a laptop CPU for the 50-pc configuration).
